\section{Completions, bounded charts, and cardinality heights}

\subsection{Exact completion of the stable choices}
\label{sec:stable-completion}
Let a finite digraph have an involution $\theta$ on vertices reversing
arcs. An individually fixed arc has the form $e:u\to\theta u$.
Let $D$ be an RC-invariant DAG from which all such candidate arcs
have been removed. Write $a_D(v)$ and $b_D(v)$ for its longest
incoming and outgoing path lengths and $L_D$ for its longest path.
Reachability below includes paths of length zero.

\begin{lemma}[At most one fixed arc on a path]
\label{lem:one-fixed-arc}
An RC-invariant DAG has no directed path containing two
individually fixed arcs.
\end{lemma}
\begin{proof}
If $e:u\to\theta u$ and $f:v\to\theta v$ occur in that order,
let $P\colon\theta u\leadsto v$ be the intervening path. Then
$e,P,f,\RC(P)$ is a nonempty closed directed walk, a contradiction.
The argument includes an empty $P$.
\end{proof}

\begin{theorem}[Stable completion by 2-SAT]
\label{thm:stable-completion}
For a subset $R$ of the removed fixed arcs, $D\cup R$ is acyclic
if and only if it satisfies the following prohibitions:
\begin{enumerate}
\item Do not restore $e:u\to\theta u$ if $D$ has a path
      $\theta u\leadsto u$.
\item Do not restore distinct $e:u\to\theta u$ and
      $f:v\to\theta v$ together if $D$ has a path
      $\theta u\leadsto v$.
\end{enumerate}
For every acyclic completion,
\begin{equation}
 L(D\cup R)=\max\{L_D,\ \max_{e:u\to\theta u\in R}(2a_D(u)+1)\}.
 \label{eq:stable-length}
\end{equation}
Thus, after fixing the nonstable PCR selections in an even-order
RC de Bruijn graph, existence of an MDS completion with $L\le\ell$
is decided by 2-SAT in polynomial time in the explicit graph size.
\end{theorem}
\begin{proof}
Each prohibition is necessary by closing one fixed arc or applying
\Cref{lem:one-fixed-arc}. Conversely a cycle must use a restored
arc. If it uses exactly one it violates the first prohibition;
otherwise two consecutive restored arcs on the cycle violate the
second. For the length formula, every path uses at most one such
arc. The longest prefix ending at $u$ and suffix starting at
$\theta u$ can be concatenated with it: any repeated vertex would
give a cycle. RC gives $b_D(\theta u)=a_D(u)$, proving equality.

For MDS completion delete the fixed nonstable selections and all
individually fixed arcs to form $D$. If $D$ is cyclic or $L_D>\ell$,
reject. Otherwise introduce a variable for restoring each fixed
arc, require exactly one per stable pair, and add the two types of
prohibition as unary and binary clauses. Also forbid any restored
arc with $2a_D(u)+1>\ell$. These are exactly the necessary and
sufficient constraints. Reachability and DAG distances are polynomial.
The implication graph decides 2-SAT by checking that no variable
and its negation belong to one strongly connected component.
\end{proof}
The implication-graph algorithm is classical
\cite{AspvallPlassTarjan1979}; the complete RC reachability and
length reductions are proved above. The algorithm fixes the
nonstable selections first; it is not a polynomial algorithm for
the global MDS optimization problem or for counting all 2-SAT models.
For $\ell=106$, a restored fixed arc must have $a_D(u)\le52$.

\begin{proposition}[The complete stable fibre of $M_{107}$]
\label{prop:stable-fibre}
Keeping all nonstable selections of the canonical $M_{107}$ fixed
leaves exactly 128 RC MDS: 121 of the 128 stable choices are forced
and seven are free. The resulting cube has ordinary-component
minimum profile
\[
 107:8,\qquad109:8,\qquad111:16,\qquad113:32,\qquad115:64.
\]
The 128 endpoints belong to distinct ordinary components. None of
these fixed-exterior completions has $L\le106$.
\end{proposition}
\begin{proof}[Exact certificates]
The stable-cube certificate reconstructs all fixed arcs, all
reachability prohibitions and all costs in \eqref{eq:stable-length}.
It verifies the 128 DAGs and, for each, a one-hit cycle one arc
longer than its length. These certify the component minima, not
just the fibre minima. Distinct stable assignments in this cube
are neither equal nor global complements, since 121 bits are
fixed. \Cref{lem:stable-pcr-parity} separates their components.
A 107-arc path using only nonstable PCRs remains after every stable
change, and independently obstructs the target 106.
\end{proof}

\subsubsection{Allowing one extra selected arc}
Let $N$ be the number of PCRs. At cost at most $N+1$, an RC
feedback set has exactly one selection in every nonstable PCR:
surplus there comes in RC pairs and costs at least two. At cost
$N+1$, exactly one stable PCR is double. Its selected pair is
either its two individually fixed arcs or a nonfixed RC pair.

\begin{corollary}[Complete fixed-exterior oracle for excess at most one]
\label{cor:excess-one-oracle}
For even $q,k$ and free alphabet complement, fix an RC selection
of one arc per nonstable PCR. Existence of an RC feedback set
extending it with cost at most $N+1$ and $L\le\ell$ is decided
by at most $1+\sum_{P\ \mathrm{stable}}|P|/2$ instances of 2-SAT.
This is polynomial in the explicit graph size. For DNA $k=8$
the number of branches is $1+128+366=495$.
\end{corollary}
\begin{proof}
The preceding classification is exhaustive because every PCR
must be hit. Take one branch with no exceptional PCR, or choose
the exceptional PCR and its double selection. A stable PCR of
period $p$ has one fixed pair and $(p-2)/2$ nonfixed pairs, giving
$p/2$ exceptional branches. Construct $D$ by deleting the fixed
nonstable selections, the chosen exceptional nonfixed pair if any,
and all individually fixed arcs. Reject cyclic or overlong $D$.
Ordinary stable pairs require exactly one restored fixed arc.
For the exceptional pair, force both restoration variables to zero
if both fixed arcs were selected, or both to one if a nonfixed
pair was selected. All constraints are 2-CNF.
\Cref{thm:stable-completion} applies to any subset of restored
fixed arcs, including these forced cases, and gives exact clauses.
Their union over the branches is precisely the desired family.
The stable PCRs have 988 arcs at DNA order eight, yielding 495.
\end{proof}
This corollary combines the completion theorem with the surplus
classification; it does not choose the 4051 nonstable RC-PCR
groups of the global problem. Its finite regression includes
satisfiable and unsatisfiable 2-SAT branches, not just cases
rejected before the SAT step.

\subsubsection{Rank projection and cuts with the correct surplus}
\begin{proposition}[Continuous ranks for Boolean endpoints]
\label{prop:boolean-ranks}
For a Boolean arc indicator $n$ and integer $\ell\ge0$, the residual
is a DAG with $L\le\ell$ if and only if there is a real rank
$r\in[0,\ell]^V$ such that
\[
 r(h(e))-r(t(e))\ge1-(\ell+1)n(e).
\]
If $n$ is RC invariant the rank may have mate sum $\ell$.
\end{proposition}
\begin{proof}
An integral longest-path rank proves necessity, including on
selected arcs where the right side is $-\ell$. Conversely summing
along zero arcs excludes cycles and paths longer than $\ell$.
Replace $r(v)$ by $(r(v)+\ell-r(\theta v))/2$ for the last claim.
\end{proof}
Integrality of the signed clock in \Cref{thm:bounded-chart} remains
essential; this proposition assumes Boolean $n$.

The rank system can be projected exactly to a boundary $B$ of
variable arcs, retaining every variable-arc inequality on its
boundary endpoints. Let $D$ be the fixed residual DAG and keep all
inequalities $r(v)-r(u)\ge d_D(u,v)$ for reachable boundary pairs,
where $d_D$ is the longest path distance. Also keep
$a_D(v)\le r(v)\le\ell-b_D(v)$ and check $L_D\le\ell$.
To extend feasible boundary values, set the rank of any vertex
$v$ to the maximum of $a_D(v)$ and all
$r(u)+d_D(u,v)$ with $u\in B$ reaching $v$.
The boundary values are preserved; the suffix bounds keep every
extended rank at most $\ell$, and all fixed arc inequalities hold.
This proves exact elimination. It may create many boundary-pair
constraints and is not an automatic reduction in total model size.

\begin{corollary}[Two-arm inequality with stable surplus]
\label{cor:surplus-two-arm}
Let a stable PCR $P$ have fixed arcs $e_i:u_i\to\theta u_i$,
$i=0,1$. For each $i$ choose a path of $h$ nonfixed arcs ending
at $u_i$, and let $U$ be the union of their arc sets. Every RC
feedback set $M$ with $L(M)\le2h$ satisfies
\[
 \sum_{e\in U}\ind_M(e)+\bigl(|M\cap P|-1\bigr)\ge1.
\]
\end{corollary}
\begin{proof}
The PCR is hit. If it is hit at least twice the inequality is
immediate. Otherwise RC forces its sole selected arc to be fixed,
so the other $e_i$ is residual. If $M$ avoided both arms, RC
would also leave both reflected arms residual. The arm ending
at $u_i$, followed by $e_i$ and its reflected arm, would be a
residual walk of $2h+1$ arcs. It is either an overlong path or
contains a cycle, both impossible.
\end{proof}
The sum uses the full arc universe, including arcs that become
selectable when surplus is allowed. The surplus belongs to the
\emph{stable} PCR $P$. At MDS cost it vanishes; at cost $N+1$
it can legitimately deactivate a clause. Omitting it is false:
for binary $k=4$, the RC feedback set
\[
 \{0000,0011,0100,0101,1000,1100,1101,1110,1111\}
\]
has $L=2$. The stable PCR $\{0011,0110,1001,1100\}$ is hit twice,
while neither one-arc arm $0001$ or $0110$ is selected. The
unlifted inequality would read $0\ge1$.

\subsection{A complete bounded chart of RC endpoints}
\label{sec:signed-chart}
Assume throughout this section that $q$ and $k\ge2$ are even,
the alphabet complement has no fixed letter, and a base indicator
$m$ is RC invariant with exactly one selection per PCR. Lower
vertices have odd length and occur in distinct pairs $\{u,\theta u\}$.
The base need not be acyclic for the algebraic representation below.

\begin{lemma}[Stable PCRs and a parity restriction]
\label{lem:stable-pcr-parity}
Every RC-stable PCR has exactly two individually RC-fixed arcs.
An RC selection of one arc in that PCR chooses one of them.
Every integral combination of paired transfers
\[
 g_{u,a}=\ind_{ua}-\ind_{au}
             +\ind_{\RC(ua)}-\ind_{\RC(au)}
\]
preserves the parity of each individually fixed arc coefficient.
If two RC MDS differ by an integral gradient, all their fixed-arc
bits are either unchanged or simultaneously complemented.
\end{lemma}
\begin{proof}
Write $\RC(w)=\rot^j w$ on an orbit of primitive period $p$.
The induced reflection of letter positions cannot fix a position,
since a letter is different from its complement. Thus $p$ is even
and $j$ is even. Fixed rotations solve $2i=j\pmod p$, giving two
solutions. A one-element RC-invariant selection must be fixed.
At a fixed arc the two summands in each RC pair in $g_{u,a}$
coincide, so its coefficient is even.
For $n-m=\nabla f$, comparison of mate arcs makes
$f(v)+f(\theta v)=C$ constant. On $e:u\to\theta u$,
$n(e)-m(e)=2f(u)-C\equiv C\pmod2$; Boolean endpoints prove the claim.
\end{proof}
For DNA $k=8$, the 128 stable choices, modulo simultaneous
complementation, yield a 127-bit ambient label for ordinary
components. This is neither a complete invariant nor a claim that
all labels are realized. Omitting independent stable corrections
therefore imposes a genuine restriction on an integral transfer search.

\begin{lemma}[Reconstruction from a signed cochain]
\label{lem:signed-reconstruction}
Let $x:E\to\Z$ be RC invariant and sum to one on every PCR.
If $T\ge1$ and an integral $H$ satisfy
\begin{equation}
 H(h(e))-H(t(e))+Tx(e)\ge1,\qquad
 H(v)+H(\theta v)=T-1,
 \label{eq:signed-clock}
\end{equation}
then
\[
 n(e)=x(e)-\floor{H(t(e))/T}+\floor{H(h(e))/T}
\]
is an RC MDS indicator and $L(n)\le T-1$.
\end{lemma}
\begin{proof}
Write $H=Ta+r$, $0\le r<T$. Then
$Tn(e)+r(h(e))-r(t(e))\ge1$, so integrality gives $n(e)\ge0$.
Its sum on each PCR is one by telescoping. Since the PCRs partition
the arcs, these conditions force $n(e)\in\{0,1\}$. Rank $r$
increases along each zero arc. The mate sum gives
$r(v)+r(\theta v)=T-1$ and $a(v)+a(\theta v)=0$, making
$n=x-\nabla a$ RC invariant. Its cardinality attains the PCR lower
bound. This includes $T=1$.
\end{proof}

Choose one representative of each lower-vertex pair and a reference
letter $a_*$. In each stable PCR $P$, let $f_P^0$ be the selected
fixed arc of $m$ and $f_P^1$ the other, and set
$\eta_P=\ind_{f_P^1}-\ind_{f_P^0}$. Put
\[
 x=m+\sum_{u,\,a\ne a_*}z_{u,a}g_{u,a}
                         +\sum_{P\ \mathrm{stable}}\gamma_P\eta_P.
\]
\begin{theorem}[Complete bounded endpoint chart]
\label{thm:bounded-chart}
An RC MDS of residual length at most $T-1$ exists if and only if
\eqref{eq:signed-clock} has a solution with
\[
 z_{u,a}\in\{-2,-1,0,1,2\},\quad \gamma_P\in\{0,1\},\quad
 -T\le H(v)\le2T-1,
\]
all coordinates integral. There is no restriction on the number
of nonzero sites $u$.
\end{theorem}
\begin{proof}
Each generator preserves RC and each PCR sum, so sufficiency is
\Cref{lem:signed-reconstruction}. For necessity, take a required
endpoint $n$. In one PCR of each nonstable RC pair, transport the
selection of $m$ to that of $n$ along a rotation path shorter than
a full turn, and add its RC companion. After canonicalizing $u$,
this gives transfer coefficients $t_{u,a}\in\{-1,0,1\}$: a
coordinate can arise only from a rotation step or its reversed RC
companion in the other PCR, and different PCR pairs cannot share it.
Treat stable PCRs directly with $\gamma_P\eta_P$.

The identities
\[
 g_{\theta u,\bar a}=-g_{u,a},\qquad
 \sum_a g_{u,a}=\nabla(\ind_u-\ind_{\theta u})
\]
allow elimination of the reference coordinate. Set $b_u=t_{u,a_*}$
and $z_{u,a}=t_{u,a}-b_u$. Then $n=x+\nabla f$, where
$f=\sum_u b_u(\ind_u-\ind_{\theta u})$ takes values in
$\{-1,0,1\}$ and is antisymmetric. Thus $|z_{u,a}|\le2$.
The bounded autodual-rank proof in \Cref{lem:autodual-rank} applies
also here at $k=2$: take an integral rank $r\in[0,T-1]$ for the
residual of $n$, with mate sum $T-1$. Setting $H=r-Tf$ gives the
stated bounds and all inequalities.
\end{proof}
At DNA order eight the chart uses $8192\cdot3=24576$ transfer
coordinates, 128 stable bits and 8192 independent clock coordinates,
for 32,896 integer variables. A signed $x$ need not be Boolean or
acyclic. The theorem concerns its reconstructed endpoint, not the
legality of intermediate local moves. The three-site exclusion below
allows arbitrary integer coefficients on those sites and arbitrary
retiming; its local scope is not a limitation of this complete chart.

\subsection{Bounded heights and cardinality above the PCR minimum}
\label{sec:height-excess}
The MDS objective fixes cardinality first. A complementary question fixes
the path bound and minimizes cardinality. Write
\[
 \tau_{\RC}(\ell)=\min\{|M|:M\text{ is an RC feedback set and }L(M)\le\ell\}.
\]
For DNA $k=8$, every such set has at least 8230 arcs. A set of
8232 arcs is not an MDS, even when it is inclusion-minimal.

\begin{theorem}[Exact antisymmetric height formulation]
\label{thm:height-formulation}
Let $G$ be a finite digraph with a vertex involution $\theta$ reversing
arcs, and let $H\ge0$ be an integer. The minimum cardinality of an
RC feedback set with residual length at most $2H$ is
\[
 \min_{\substack{h:V\to\Z,\ |h(v)|\le H\\h(\theta v)=-h(v)}}
       |\{e:u\to v:h(v)\le h(u)\}|.
\]
The statement allows fixed vertices of $\theta$.
\end{theorem}
\begin{proof}
For a feasible height $h$, retain only strictly increasing arcs.
This is an RC DAG of path length at most $2H$, proving one
inequality. Conversely take a feasible feedback set and its
residual longest incoming rank $d$. The half-integral height
$c(v)=(d(v)-d(\theta v))/2$ lies in $[-H,H]$, is antisymmetric,
and increases by at least one on residual arcs. Round nonintegral
values away from zero, keeping integers fixed. For endpoints on
one side of zero their rounding increments agree if both are
half-integral; across zero the gap grows. With exactly one
half-integral endpoint an integral lower gap bound of one means
the original gap is at least $3/2$, so losing at most $1/2$
is harmless. Thus the resulting integral $h$ still increases
on all residual arcs. It stays in $[-H,H]$ and is zero on fixed
vertices. Its nonincreasing arcs form a subset of the given
feedback set, proving the reverse inequality.
\end{proof}
Without antisymmetry, maximizing the increasing arcs is the
classical Max-$k$-Ordering formulation, here with $2H+1$ labels
\cite{KenkreEtAl2015}. The statement above supplies the exact
RC restriction for an even path bound, with a proof rather than
a continuous-relaxation assumption.

\begin{proposition}[Cost 8232 at length 105]
\label{prop:8232}
There is an RC feedback set $S$ with $|S|=8232$ and $L(S)=105$.
It is inclusion-minimal even among feedback sets without RC.
It hits 8228 PCRs once and two nonstable companion PCRs twice.
In particular
\[
 8230\le\tau_{\RC}(106)\le8232.
\]
The same upper bound holds for $\tau_{\RC}(105)$.
\end{proposition}
\begin{proof}[Exact certificates]
The word list \pathfile{r11/certificates/S8232.words.txt}, height
certificate, rank and attaining path verify the first two numbers.
The checker reconstructs the PCR partition and the stated
multiplicities. For each singly hit PCR, its rotation cycle
certifies necessity of its selected arc. For the four selected
arcs in the double PCRs, four explicitly stored one-hit cycles
of lengths six or nine provide the missing necessity witnesses.
Removing any selected arc therefore restores a directed cycle.
The PCR count proves the lower bound 8230 independently of this
construction. Neither inclusion-minimality nor these short cycles
prove that 8232 is a minimum at the required length.
\end{proof}

There are two further cost-8232 witnesses $S_A,S_B$ of length 105
in \pathfile{r13/certificates/}. They arise by adding the RC pair
$\{\texttt{GAAAAAAA},\texttt{TTTTTTTC}\}$ to the two length-107
MDS \pathfile{M107_panel_02.words.txt} and
\pathfile{M107_panel_03.words.txt}. Those MDS are in one ordinary
component and opposite sectors. In this component, two is the
minimum number of RC additions needed to reach length at most 106:
the supplied simple 108-cycle has one selected arc and no
individually fixed arcs. Any single RC addition would have to be
fixed, so it cannot repair that cycle. Gradients preserve its
intersection count. The two displayed additions attain length 105.
For the original canonical $M_{107}$, a separate exact census
instead gives minimum RC augmentation cost four. These are
fixed-base or fixed-component augmentation statements, not
global cardinality lower bounds.

Adding either word of the displayed pair alone gives a set of
8231 arcs with $L=106$, but that set is not RC invariant. This
separates an ordinary construction from the still-open RC target.

\begin{proposition}[Affine pruning is exhausted for three bases]
\label{prop:affine-pruning}
For each of $S,S_A,S_B$, no smaller Boolean RC feedback indicator
of the form
\[
 n=\ind_S-a+\nabla f,\qquad
 a:E\to\Z_{\ge0},\quad a\ne0,\quad a\text{ RC invariant},
 \quad f:V\to\R,
\]
has residual length at most 106 (replace $S$ by the chosen base).
For the first base no smaller feedback endpoint of this form
exists at any length. For $S_A,S_B$ the smallest attainable
length in this family is 107.
\end{proposition}
\begin{proof}[Reduction and exact cycle certificates]
On each PCR, gradients telescope and a feedback indicator has
positive integral sum. Hence $a$ vanishes on every singly hit
PCR, including the stable ones. Each base has just one doubled
nonstable PCR pair. In each of its two PCRs, $a$ has sum at most
one. RC and $a\ne0$ force the sums to be one, so $a$ is one
RC pair of arcs chosen from this orbit pair. There are exactly
eight possibilities, including arcs not selected in the base.
The whole de Bruijn graph is balanced, so $\sum_e\nabla f(e)=0$;
any endpoint then has cardinality 8230.

For the first base all eight choices have a certified closed
walk of nonpositive $\ind_S-a$ weight. Its weight cannot change
under a gradient; a Boolean feedback indicator must have positive
weight on every nonempty closed walk. For each of the other two
bases there is one negative and six zero-weight obstructions.
The final case has a 108-cycle of weight one, forcing $L\ge107$.
In that case subtracting the pair actually gives the corresponding
length-107 MDS, so the bound is attained. All 24 walks and their
integer weights are checked in
\pathfile{r13/certificates/affine_pruning.json}.
\end{proof}
Allowing real $f$ does not weaken the cycle obstructions. Between
Boolean endpoints an integral gradient also has an integral
potential after a constant normalization, by propagation along
the connected graph. This exhaustion concerns subtraction and
retiming. It does not allow arbitrary new inserted arcs, which
are a separate research direction.
